% Content moved from project thesis background §2.2 (Physics Simulation)
% to master thesis background §2.2, because cloth dynamics are the focus
% of the master thesis, not the project thesis.

%===================================================================%
\section{Cloth and Deformable Object Simulation}\label{sec:cloth_sim}
%===================================================================%

Deformable object simulation relaxes the rigid-body assumption by allowing the object's
shape to change under load. This introduces two coupled complexities: the state dimension
increases dramatically (many vertices/particles instead of a single rigid transform), and
contacts become configuration-dependent because deformation continuously alters the set of
potential contact points and normals~\cite{Nealen2006DeformableModels,Longhini2025ReviewClothManipulation}.
For textiles, this coupling is particularly pronounced: applied forces not only translate the
object but also create folds, wrinkles, and contact-rich self-interactions that alter future
dynamics. These properties are repeatedly identified as key contributors to the reality gap in
robotic cloth manipulation~\cite{Longhini2025ReviewClothManipulation}.

\paragraph{Common modelling approaches.}
Three modelling families are most relevant in robotics simulation.
First, mass--spring(-damper) models discretize cloth as a mesh whose vertices are connected
by springs encoding stretch and (optionally) bending penalties; such models can be made
stable with implicit integration techniques but require careful parameterization and can be
numerically stiff~\cite{BaraffWitkin1998LargeStepsCloth}.
Second, particle-based constraint methods represent cloth by particles and enforce desired
behavior via geometric (holonomic) constraints, which can be solved iteratively with high
robustness under collisions; \ac{PBD} and its variants are canonical examples~\cite{muller2007pbd,macklin2016xpbd}.
Third, continuum formulations discretize elastic energy and constitutive laws (e.g., in
\ac{FEM}) or use hybrid particle--grid methods (e.g., \ac{MPM}), providing a pathway to
parameters with clearer physical interpretation at increased computational cost~\cite{SifakisBarbic2012FEM,Jiang2016MPM}.

\paragraph{Challenges specific to cloth dynamics.}
Cloth exhibits high effective \ac{DoF}, frequent self-collision, and strong sensitivity to
material parameters such as bending stiffness, friction, and damping. Robust self-contact
and friction handling is widely recognized as a hard requirement for stable cloth simulation,
particularly under large deformations and repeated collisions~\cite{Bridson2002ClothCollisions,Longhini2025ReviewClothManipulation}.
From a learning perspective, this sensitivity implies that small modelling errors (e.g., in
friction or damping) can change grasp success rates, sliding behavior, and fold formation,
which in turn can qualitatively alter the learned policy's strategy and its transferability to
hardware~\cite{Salvato2021RealityGapSurvey,Longhini2025ReviewClothManipulation}.

\subsection{\acs{PBD} --- PhysX Particle Cloth}

PhysX provides particle-based simulation through \texttt{PxPBDParticleSystem}, whose solver
is explicitly described as using \ac{PBD} for particle--particle dynamics and as supporting
phenomena including cloth, inflatables, and other deformable behaviors~\cite{physx_particle_system_513,physx_pbd_particlesystem_api}.
In Omniverse, particle simulation is \ac{GPU}-accelerated and is documented as supporting
fluids, granular media, and cloth, including interaction with rigid bodies and
articulations~\cite{omni_physics_particles}.
At the feature level, particle cloth is treated as a particle-based \ac{PBD} mass--spring-style
system and is marked as deprecated in favor of newer deformable-body features, which
highlights that ``cloth'' in Isaac~Sim can refer to multiple, evolving solver families rather
than a single canonical model~\cite{omni_physics_particles,isaac_sim_physics_resources_510}.

\paragraph{Time-step structure and iterative projection.}
In standard \ac{PBD}, the deformable is represented by particle positions
$\mathbf{x}_i \in \mathbb{R}^3$ with masses $m_i$; desired behavior is enforced by a set of
constraints $C_j(\mathbf{x})$ (equalities or inequalities), typically solved by iterative
projection in a Gauss--Seidel fashion~\cite{muller2007pbd}.
A schematic time step consists of: (i) updating velocities under external forces,
(ii) predicting positions by an explicit step, (iii) iteratively projecting predicted positions
to reduce constraint violations, and (iv) updating velocities from the corrected positions
(and applying damping if configured)~\cite{muller2007pbd}.
For a constraint $C_j(\mathbf{x}) = 0$, linearization yields a correction direction based on the
constraint gradient. Using the inverse mass matrix $\mathbf{M}^{-1}$ (diagonal in the particle
basis), one common update can be written as~\cite{muller2007pbd,macklin2016xpbd}:
\begin{equation}
\Delta \mathbf{x} = k_j \cdot s_j \cdot \mathbf{M}^{-1}\nabla C_j(\mathbf{x}), \qquad
s_j = \frac{-C_j(\mathbf{x})}{\nabla C_j(\mathbf{x})^\top \mathbf{M}^{-1} \nabla C_j(\mathbf{x})},
\label{eq:pbd_correction}
\end{equation}
where $k_j \in [0,1]$ is a user-defined stiffness-like scaling applied per constraint.

\paragraph{Effective stiffness and dependence on $h$ and iteration count.}
A central practical implication of the iterative projection view is that the \emph{effective}
stiffness of constraints depends on numerical settings. Increasing the number of projection
iterations typically yields stiffer constraint satisfaction at fixed $h$, while increasing $h$
tends to reduce stability and can require either more iterations or smaller substeps to
maintain comparable behavior~\cite{muller2007pbd,macklin2016xpbd,Macklin2019SmallSteps}.
This is beneficial for robustness and controllability, but it complicates mapping simulator
parameters to physically measured material properties (e.g., bending stiffness in SI units),
which is a recurring issue when calibrating cloth simulation for sim-to-real
transfer~\cite{macklin2016xpbd,Salvato2021RealityGapSurvey}.

\paragraph{\acs{XPBD} compliance formulation.}
\ac{XPBD} extends \ac{PBD} by introducing a compliance parameter and an explicit Lagrange
multiplier state to reduce time-step dependence and recover meaningful constraint force
estimates~\cite{macklin2016xpbd}.
For compliance $\alpha \ge 0$ and time step $h$, the incremental multiplier update for a
constraint can be expressed as~\cite{macklin2016xpbd}:
\begin{equation}
\Delta \lambda =
\frac{-C(\mathbf{x}) - \frac{\alpha}{h^2}\lambda}{\sum_i w_i \|\nabla_{\mathbf{x}_i} C(\mathbf{x})\|^2 + \frac{\alpha}{h^2}},
\qquad
\Delta \mathbf{x}_i = w_i \nabla_{\mathbf{x}_i} C(\mathbf{x}) \Delta \lambda,
\label{eq:xpbd_update}
\end{equation}
with inverse masses $w_i = 1/m_i$.
This formulation decouples the user parameter $\alpha$ from the iteration count in a way
that improves interpretability of stiffness-like behavior across different $h$ and solver
settings~\cite{macklin2016xpbd}.

\paragraph{Why \ac{PBD} is used in PhysX for particle cloth.}
\ac{PBD} is well-aligned with \ac{GPU} execution because constraints can be partitioned into
independent sets and processed in parallel; graph coloring is a common strategy to schedule
constraint batches without write conflicts on shared particles~\cite{FratarcangeliPellacini2013GPUPBD}.
Omniverse documentation explicitly emphasizes \ac{GPU}-accelerated \ac{PBD} particles and
their interaction with rigid bodies and articulations, which is essential for cloth manipulation
scenes in robotics~\cite{omni_physics_particles}.
In addition, the PhysX particle system API is designed around a dedicated \ac{PBD} solver
type (\texttt{PxPBDParticleSystem}), reinforcing that particle cloth is treated as a specific,
robust, high-throughput model family within the PhysX multiphysics portfolio~\cite{physx_particle_system_513,physx_pbd_particlesystem_api}.

\subsection{\acs{VBD} --- Newton Thin Deformables}

While PhysX particle cloth is based on explicit prediction and iterative position projection,
implicit methods aim to improve stability under stiffness and large time steps by solving a
time-discrete optimization problem. \ac{VBD} is introduced as a block coordinate descent
solver for the variational form of implicit Euler time integration, targeting convergence to
the implicit Euler solution while maintaining unconditional stability under limited iteration
budgets~\cite{Chen2024VBD}.

\paragraph{Variational implicit Euler and global energy minimization.}
Let $\mathbf{x}_t$ denote positions, $\mathbf{v}_t$ velocities, $h$ the time step, and
$\mathbf{M}$ the mass matrix. In the variational formulation, the next positions are obtained
by minimizing a global energy that combines inertia and potential energy~\cite{Chen2024VBD}:
\begin{equation}
\mathbf{x}_{t+1} = \arg\min_{\mathbf{x}}\, G(\mathbf{x}), \qquad
G(\mathbf{x}) = \frac{1}{2h^2}\|\mathbf{x}-\mathbf{y}\|_{\mathbf{M}}^2 + E(\mathbf{x}),
\label{eq:vbd_global_energy}
\end{equation}
where $\mathbf{y}=\mathbf{x}_t+h\mathbf{v}_t+h^2\mathbf{a}_\mathrm{ext}$ and
$E(\mathbf{x})$ is the total potential energy.
Velocities follow from the implicit Euler relationship
$\mathbf{v}_{t+1}=(\mathbf{x}_{t+1}-\mathbf{x}_t)/h$~\cite{Chen2024VBD}.

\paragraph{Vertex-level block coordinate descent and local Newton step.}
\ac{VBD} solves the global problem by iterating over vertices and updating one vertex at a
time while temporarily holding the others fixed, yielding a vertex-level Gauss--Seidel
pattern with favorable parallelization via vertex coloring~\cite{Chen2024VBD}.
For vertex $i$, a local energy is minimized:
\begin{equation}
\mathbf{x}_i \leftarrow \arg\min_{\mathbf{x}_i} G_i(\mathbf{x}), \qquad
G_i(\mathbf{x})=\frac{m_i}{2h^2}\|\mathbf{x}_i-\mathbf{y}_i\|^2+\sum_{j\in F_i}E_j(\mathbf{x}),
\label{eq:vbd_local_energy}
\end{equation}
where $F_i$ is the set of force/energy elements incident to vertex $i$.
Because $G_i$ has only three degrees of freedom, a local Newton step is formed by solving
a $3\times 3$ linear system per vertex~\cite{Chen2024VBD}:
\begin{equation}
\mathbf{H}_i \Delta \mathbf{x}_i = \mathbf{f}_i,
\label{eq:vbd_newton_step}
\end{equation}
with $\mathbf{H}_i$ the local Hessian and $\mathbf{f}_i$ the local gradient/force term.

\paragraph{Advantages for stiff, contact-rich thin deformables.}
The primary advantages highlighted for \ac{VBD} are: (i) unconditional stability inherited
from the energy descent property, (ii) convergence toward the implicit Euler solution as
iterations increase, and (iii) \ac{GPU}-oriented parallelization via vertex coloring with small,
localized linear solves~\cite{Chen2024VBD}.
These properties are attractive for cloth-like thin deformables, which often exhibit
near-inextensible stretch (high stiffness) combined with softer bending and frequent
contact events; implicit/variational formulations can tolerate such stiffness ratios at larger
time steps than explicit solvers with comparable iteration budgets~\cite{Chen2024VBD}.

\paragraph{Newton in the robotics simulation ecosystem.}
Newton is presented as an open-source, \ac{GPU}-accelerated physics engine designed for
robotics simulation and robot learning, built on NVIDIA Warp and intended to integrate with
robot learning frameworks such as IsaacLab~\cite{newton_physics_doc,nvidia_newton_announcement_2025}.
In NVIDIA's robotics communication, Newton is explicitly associated with thin-deformable
simulation via a \ac{VBD} solver and with multiphysics capabilities that include implicit
\ac{MPM} and other solver components, motivated by stability and performance for learning
workflows~\cite{nvidia_isaaclab_newton_blog_2025,nvidia_newton_announcement_2025}.
Accordingly, \ac{VBD} is positioned as a solver choice when increased stability under stiff
cloth behavior and larger time steps is required, complementing the high-throughput
\ac{PBD}-based particle cloth pathway in PhysX-centric stacks~\cite{nvidia_isaaclab_newton_blog_2025,Chen2024VBD}.

\subsection{Comparison: \acs{PBD} vs.\ \acs{VBD} for Cloth in \ac{RL} Workflows}

\paragraph{Mathematical targets and implications.}
The core distinction is the optimization target.
\ac{PBD} advances dynamics by predicting positions and then \emph{projecting} them to reduce
constraint violations, where the resulting behavior depends on solver scheduling and
iteration count, and where constraint ``stiffness'' is typically not time-step invariant without
extensions~\cite{muller2007pbd,macklin2016xpbd}.
\ac{VBD}, in contrast, directly minimizes the variational implicit Euler energy and therefore
targets a principled implicit integration objective; stability is maintained by construction via
energy descent, and increased iteration budgets refine the approximation toward the implicit
solution~\cite{Chen2024VBD}.

\paragraph{Practical tradeoffs for learning-centric simulation.}
For \ac{RL}, \ac{PBD}-based particle cloth is attractive because it is simple, robust under
frequent contacts, and aligned with \ac{GPU}-parallel execution through constraint batching
and coloring strategies; this supports the high environment throughput required for policy
optimization~\cite{omni_physics_particles,FratarcangeliPellacini2013GPUPBD,Makoviychuk2021IsaacGym}.
However, because effective stiffness depends on $h$ and iteration count, parameter tuning
can be difficult to interpret physically and can amplify sim-to-real sensitivity if policies rely
on artefacts (e.g., overly damped folds or under-resolved frictional sticking)~\cite{macklin2016xpbd,Salvato2021RealityGapSurvey}.
\ac{VBD} is attractive when stability under stiffness, large time steps, and mixed constraint
types is prioritized, particularly for systems that combine near-inextensible stretch with
softer bending; such regimes are common in cloth manipulation and can be challenging for
explicit projection methods at fixed compute budgets~\cite{Chen2024VBD,nvidia_isaaclab_newton_blog_2025}.

\paragraph{Relation to the Omniverse ecosystem and thesis choice.}
In Omniverse/Isaac~Sim, particle cloth is documented as deprecated in favor of newer
deformable-body features, indicating an ongoing transition toward more general deformable
schemas and alternative solver families~\cite{omni_physics_particles,isaac_sim_physics_resources_510}.
For the experiments in this thesis, the PhysX/\ac{PBD} particle cloth pathway is motivated
by \ac{GPU} throughput and robustness for large-scale training in Isaac~Sim~5.1, whereas
Newton/\ac{VBD} is conceptually aligned with scenarios where greater stability under stiff
cloth behavior and larger time steps is required (e.g., when cloth parameters or contact
conditions become numerically challenging for projection-based solvers)~\cite{omni_physics_particles,Chen2024VBD,nvidia_isaaclab_newton_blog_2025}.
% TODO: verify which cloth/deformable features are enabled in the exact Isaac~Sim~5.1.0 installation used, and document the final solver selection in the Methodology chapter.

\paragraph{Rationale for an integrated cloth--rigid simulation stack.}
% Content moved verbatim from project thesis background §2.3 (Isaac Sim and
% IsaacLab Architecture, Rationale subsection), because it justifies the
% simulator choice specifically from the cloth-dynamics perspective.
Integrated cloth and deformable dynamics are supported within the Isaac Sim / PhysX stack through GPU-accelerated particle-based simulation that can represent cloth-like behavior and interact with articulations and other physics objects within a shared simulation scene~\cite{omni_physics_particles}. In addition, IsaacLab explicitly highlights support for deformable objects (including cloth) within a framework built on Isaac Sim and designed for multi-modal learning, indicating that deformable interactions are treated as a first-class concern in the simulation-to-learning pipeline~\cite{Mittal2025IsaacLab}. This integration is beneficial because it avoids ad-hoc coupling between an external deformable simulator and a separate rigid-body engine, which would otherwise require custom synchronization, contact handling, and state transfer between solvers.

% ---------------------------------------------------------------------------
% BibTeX keys used in this section (cloth/deformable):
%  - physx_particle_system_513              (PhysX particle system: PBD solver, cloth support)
%  - physx_pbd_particlesystem_api           (PhysX PxPBDParticleSystem API reference)
%  - omni_physics_particles                 (Omni Physics particles: GPU PBD particles, cloth, deprecation)
%  - isaac_sim_physics_resources_510        (Isaac Sim physics limitations: particle cloth deprecation)
%  - Makoviychuk2021IsaacGym                (NeurIPS 2021 Isaac Gym: GPU simulation for RL)
%  - Salvato2021RealityGapSurvey            (IEEE Access survey on sim-to-real in RL)
%  - Longhini2025ReviewClothManipulation    (Annual Review of robotic cloth manipulation)
%  - Nealen2006DeformableModels             (STAR: deformable model families and tradeoffs)
%  - BaraffWitkin1998LargeStepsCloth        (SIGGRAPH: stable cloth simulation with larger steps)
%  - Bridson2002ClothCollisions             (SIGGRAPH: robust cloth collision/contact/friction)
%  - SifakisBarbic2012FEM                   (SIGGRAPH course: FEM simulation fundamentals)
%  - Jiang2016MPM                           (SIGGRAPH/TOG: MPM for continuum materials)
%  - muller2007pbd                          (JVCIR: Position Based Dynamics)
%  - macklin2016xpbd                        (MIG: XPBD compliance and Lagrange multipliers)
%  - FratarcangeliPellacini2013GPUPBD       (JGT: GPU PBD with graph coloring)
%  - Macklin2019SmallSteps                  (Small Steps in Physics Simulation: substepping insight)
%  - Chen2024VBD                            (TOG: Vertex Block Descent)
%  - newton_physics_doc                     (Newton Physics documentation)
%  - nvidia_newton_announcement_2025        (NVIDIA blog: announcing Newton)
%  - nvidia_isaaclab_newton_blog_2025       (NVIDIA blog: Isaac Lab + Newton cloth manipulation)
% ---------------------------------------------------------------------------

% ---------------------------------------------------------------------------
% BibTeX entries — ensure these are in shared/bibliography/literature.bib:
%
% @online{physx_particle_system_513,
%   author  = {{NVIDIA}},
%   title   = {Particle System --- PhysX 5.1.3 Documentation},
%   year    = {2023},
%   url     = {https://nvidia-omniverse.github.io/PhysX/physx/5.1.3/docs/ParticleSystem.html},
%   urldate = {2026-03-03},
% }
%
% @online{physx_pbd_particlesystem_api,
%   author  = {{NVIDIA}},
%   title   = {{PxPBDParticleSystem} --- PhysX API Reference},
%   year    = {2023},
%   url     = {https://nvidia-omniverse.github.io/PhysX/physx/5.3.1/_api_build/class_px_p_b_d_particle_system.html},
%   urldate = {2026-03-03},
% }
%
% @online{omni_physics_particles,
%   author  = {{NVIDIA}},
%   title   = {Particles --- Omni Physics Documentation},
%   year    = {2025},
%   url     = {https://docs.omniverse.nvidia.com/kit/docs/omni_physics/latest/dev_guide/particles/particles.html},
%   urldate = {2026-03-03},
% }
%
% @article{Longhini2025ReviewClothManipulation,
%   author  = {Longhini, Alberta and Wang, Yufei and Garcia-Camacho, Irene and Blanco-Mulero, David and Moletta, Marco and Welle, Michael and Aleny{\`a}, Guillem and Yin, Hang and Erickson, Zackory and Held, David and Borr{\`a}s, J{\'u}lia and Kragic, Danica},
%   title   = {A Review of Robotic Cloth Manipulation},
%   journal = {Annual Review of Control, Robotics, and Autonomous Systems},
%   year    = {2025},
%   doi     = {10.1146/annurev-control-022723-033252},
% }
%
% @article{Nealen2006DeformableModels,
%   author  = {Nealen, Andrew and M{\"u}ller, Matthias and Keiser, Richard and Boxerman, Eddy and Carlson, Mark},
%   title   = {Physically Based Deformable Models in Computer Graphics},
%   journal = {Computer Graphics Forum},
%   volume  = {25},
%   number  = {4},
%   pages   = {809--836},
%   year    = {2006},
%   doi     = {10.1111/j.1467-8659.2006.01000.x},
% }
%
% @inproceedings{BaraffWitkin1998LargeStepsCloth,
%   author    = {Baraff, David and Witkin, Andrew},
%   title     = {Large Steps in Cloth Simulation},
%   booktitle = {Proceedings of the 25th Annual Conference on Computer Graphics and Interactive Techniques (SIGGRAPH '98)},
%   year      = {1998},
%   pages     = {43--54},
%   doi       = {10.1145/280814.280821},
% }
%
% @inproceedings{Bridson2002ClothCollisions,
%   author    = {Bridson, Robert and Fedkiw, Ronald and Anderson, John},
%   title     = {Robust Treatment of Collisions, Contact and Friction for Cloth Animation},
%   booktitle = {Proceedings of the 29th Annual Conference on Computer Graphics and Interactive Techniques (SIGGRAPH '02)},
%   year      = {2002},
%   pages     = {594--603},
%   doi       = {10.1145/566570.566623},
% }
%
% @inproceedings{SifakisBarbic2012FEM,
%   author    = {Sifakis, Eftychios and Barbi{\v{c}}, Jernej},
%   title     = {{FEM} Simulation of 3D Deformable Solids: A Practitioner's Guide to Theory, Discretization and Model Reduction},
%   booktitle = {ACM SIGGRAPH 2012 Courses},
%   year      = {2012},
%   doi       = {10.1145/2343483.2343501},
% }
%
% @article{Jiang2016MPM,
%   author  = {Jiang, Chenfanfu and Schroeder, Craig and Teran, Joseph and Stomakhin, Alexey and Selle, Andrew},
%   title   = {The Material Point Method for Simulating Continuum Materials},
%   journal = {ACM Transactions on Graphics},
%   volume  = {35},
%   number  = {4},
%   year    = {2016},
%   doi     = {10.1145/2897826.2927348},
% }
%
% @article{muller2007pbd,
%   author  = {M{\"u}ller, Matthias and Heidelberger, Bruno and Hennix, Marcus and Ratcliff, John},
%   title   = {Position Based Dynamics},
%   journal = {Journal of Visual Communication and Image Representation},
%   volume  = {18},
%   number  = {2},
%   pages   = {109--118},
%   year    = {2007},
%   doi     = {10.1016/j.jvcir.2007.01.005},
% }
%
% @inproceedings{macklin2016xpbd,
%   author    = {Macklin, Miles and M{\"u}ller, Matthias and Chentanez, Nuttapong},
%   title     = {XPBD: Position-Based Simulation of Compliant Constrained Dynamics},
%   booktitle = {Proceedings of the 9th International Conference on Motion in Games (MIG '16)},
%   year      = {2016},
%   pages     = {49--54},
%   doi       = {10.1145/2994258.2994272},
% }
%
% @article{FratarcangeliPellacini2013GPUPBD,
%   author  = {Fratarcangeli, Marco and Pellacini, Fabio},
%   title   = {A GPU-Based Implementation of Position Based Dynamics for Interactive Deformable Bodies},
%   journal = {Journal of Graphics Tools},
%   volume  = {17},
%   number  = {3},
%   pages   = {59--66},
%   year    = {2013},
%   doi     = {10.1080/2165347X.2015.1030525},
% }
%
% @online{Macklin2019SmallSteps,
%   author  = {Macklin, Miles and Storey, Kier and Lu, Michelle and Terdiman, Pierre and Chentanez, Nuttapong and Jeschke, Stefan and M{\"u}ller, Matthias},
%   title   = {Small Steps in Physics Simulation},
%   year    = {2019},
%   url     = {https://mmacklin.com/smallsteps.pdf},
%   urldate = {2026-03-03},
% }
%
% @article{Chen2024VBD,
%   author  = {Chen, Anka He and Liu, Ziheng and Yang, Yin and Yuksel, Cem},
%   title   = {Vertex Block Descent},
%   journal = {ACM Transactions on Graphics},
%   volume  = {43},
%   number  = {4},
%   pages   = {1--16},
%   year    = {2024},
%   doi     = {10.1145/3658179},
% }
%
% @online{newton_physics_doc,
%   author  = {{Newton Physics Contributors}},
%   title   = {Newton Physics Documentation},
%   year    = {2025},
%   url     = {https://newton-physics.github.io/newton/},
%   urldate = {2026-03-03},
% }
%
% @online{nvidia_newton_announcement_2025,
%   author  = {{NVIDIA}},
%   title   = {Announcing Newton, an Open-Source Physics Engine for Robotics Simulation},
%   year    = {2025},
%   url     = {https://developer.nvidia.com/blog/announcing-newton-an-open-source-physics-engine-for-robotics-simulation/},
%   urldate = {2026-03-03},
% }
%
% @online{nvidia_isaaclab_newton_blog_2025,
%   author  = {{NVIDIA}},
%   title   = {Train a Quadruped Locomotion Policy and Simulate Cloth Manipulation with NVIDIA Isaac Lab and Newton},
%   year    = {2025},
%   url     = {https://developer.nvidia.com/blog/train-a-quadruped-locomotion-policy-and-simulate-cloth-manipulation-with-nvidia-isaac-lab-and-newton/},
%   urldate = {2026-03-03},
% }
% ---------------------------------------------------------------------------
